% GENERATED FILE. Edit canonical lessons, then run npm run generate:books.
% canonical-source-hash: fnv1a64:d9372fce85764581
% canonical-lessons: HI-C42-come, HI-C42-sit, HI-C42-stand, HI-C42-give

\chapter{More Everyday Verbs}
\label{ch:hi-more-verbs}
\hlchaptermodality{42}

\begin{chapteropening}
\textbf{By the end of this chapter:} \emph{I can name four everyday actions in Hindi and recognise the infinitive ending they share.}

Complete HI-C42-give and say all four verbs in order.
\end{chapteropening}

\section[\texorpdfstring{\hi{आना} / \hi{आता} / \hi{आती} … (ānā / ātā / ātī …)}{आना / आता / आती … (ānā / ātā / ātī …)}]{\hi{आना} / \hi{आता} / \hi{आती} / \hi{आते} — to come and its present forms}

\label{lesson:HI-C42-come}



\begin{quote}
A short run of everyday verbs, one at a time.
\end{quote}

\subsection*{You'll want to know: \hi{आना}}
\textbf{\hi{आना}} (\emph{ānā}) — \textquotedblleft{}to come\textquotedblright{}.

Every Hindi infinitive ends in \textbf{-\hi{ना}} (\emph{-nā}), so this is 'to come' in the form a dictionary lists.

\begin{cousinweb}
\textbf{\hi{आना}} (\emph{ānā}) is a tadbhava, a word inherited from Sanskrit and reshaped by sound change. CDIAL lists it under \emph{āpayati}, \textquotedblleft{}causes to reach\textquotedblright{}, from the verb whose old form \emph{āpa} meant \textquotedblleft{}came towards\textquotedblright{}; Prakrit already had \emph{āvaï}, \textquotedblleft{}comes\textquotedblright{} (CDIAL 1200).
\end{cousinweb}

\begin{grammarlens}[title={What you've built}]
Take off \textbf{-\hi{ना}} and add the endings you already know: \textbf{\hi{आता}} (\emph{ātā}, masculine), \textbf{\hi{आती}} (\emph{ātī}, feminine), \textbf{\hi{आते}} (\emph{āte}, respectful or plural). The action-chaining lesson will put \textbf{\hi{आता} \hi{हूँ}} after a chained action; now the form will already be yours.
\end{grammarlens}

\subsection*{Your turn}
Say these aloud:

\begin{itemize}
\raggedright
  \item \emph{ānā}
  \item \textbf{\hi{आता} / \hi{आती} / \hi{आते}} (\emph{ātā / ātī / āte})
  \item it again, to someone standing in a doorway
\end{itemize}

\begin{tcolorbox}[breakable,colback=teal!4,colframe=teal!35!black,title={Before you move on}]
What does \textbf{\hi{आना}} mean? (\textquotedblleft{}to come\textquotedblright{}.) Give its masculine, feminine, and respectful or plural forms. (\textbf{\hi{आता} / \hi{आती} / \hi{आते}}.)
\end{tcolorbox}
\section[\texorpdfstring{\hi{बैठना} (baiṭhnā)}{बैठना (baiṭhnā)}]{\hi{बैठना} (baiṭhnā) — to sit}

\label{lesson:HI-C42-sit}



\begin{quote}
Before the new one: what did \emph{ānā} mean?
\end{quote}

\subsection*{You'll want to know: \hi{बैठना}}
\textbf{\hi{बैठना}} (\emph{baiṭhnā}) — \textquotedblleft{}to sit\textquotedblright{}.

Same \emph{-\hi{ना}} ending. Once you see it, you can spot an infinitive anywhere in a text without knowing the verb.

\begin{cousinweb}
\textbf{\hi{बैठना}} (\emph{baiṭhnā}) is a tadbhava, a word inherited from Sanskrit and reshaped by sound change. It comes from \emph{upaviśati}, \textquotedblleft{}sits down\textquotedblright{} (CDIAL 2245). The \emph{-ṭh-} goes back to its past form \emph{upaviṣṭa}, \textquotedblleft{}seated\textquotedblright{}.
\end{cousinweb}

\begin{grammarlens}[title={What you've built}]
One more everyday action. Three more follow, and each reuses the ones before it.
\end{grammarlens}

\subsection*{Your turn}
Say these aloud:

\begin{itemize}
\raggedright
  \item \emph{baiṭhnā}
  \item it again, to someone standing in a doorway
  \item it after \emph{ānā}, so the two sit together
\end{itemize}

\begin{tcolorbox}[breakable,colback=teal!4,colframe=teal!35!black,title={Before you move on}]
What does \hi{बैठना} mean? (\textquotedblleft{}to sit\textquotedblright{}.) Say it once more.
\end{tcolorbox}
\section[\texorpdfstring{\hi{खड़ा} \hi{होना} (khaṛā honā)}{खड़ा होना (khaṛā honā)}]{\hi{खड़ा} \hi{होना} (khaṛā honā) — to stand}

\label{lesson:HI-C42-stand}



\begin{quote}
Before the new one: what did \emph{baiṭhnā} mean?
\end{quote}

\subsection*{You'll want to know: \hi{खड़ा} \hi{होना}}
\textbf{\hi{खड़ा} \hi{होना}} (\emph{khaṛā honā}) — \textquotedblleft{}to stand\textquotedblright{}.

Two words: an adjective \emph{khaṛā}, 'upright', plus \emph{honā}, 'to be'. Hindi builds many verbs this way, from a word plus a light verb.

\begin{cousinweb}
\textbf{\hi{खड़ा} \hi{होना}} (\emph{khaṛā honā}) is two inherited words. \textbf{\hi{खड़ा}} (\emph{khaṛā}), \textquotedblleft{}standing\textquotedblright{}, is listed by CDIAL under Sanskrit \emph{khaḍaka}, \textquotedblleft{}a bolt, a post\textquotedblright{}, for which it sets up an older sense \textquotedblleft{}erect\textquotedblright{} (CDIAL 3784). The same entry reconstructs a verb \emph{*khaḍati}, \textquotedblleft{}stands\textquotedblright{}, but lists only a few hill-dialect verbs under it; \emph{khaṛā} sits with the post. \textbf{\hi{होना}} (\emph{honā}), \textquotedblleft{}to be, to become\textquotedblright{}, comes from Sanskrit \emph{bhavati}, \textquotedblleft{}becomes, is\textquotedblright{} (CDIAL 9416).
\end{cousinweb}

\begin{grammarlens}[title={What you've built}]
One more everyday action. Three more follow, and each reuses the ones before it.
\end{grammarlens}

\subsection*{Your turn}
Say these aloud:

\begin{itemize}
\raggedright
  \item \emph{khaṛā honā}
  \item it again, to someone standing in a doorway
  \item it after \emph{baiṭhnā}, so the two sit together
\end{itemize}

\begin{tcolorbox}[breakable,colback=teal!4,colframe=teal!35!black,title={Before you move on}]
What does \hi{खड़ा} \hi{होना} mean? (\textquotedblleft{}to stand\textquotedblright{}.) Say it once more.
\end{tcolorbox}
\section[\texorpdfstring{\hi{देना} (denā)}{देना (denā)}]{\hi{देना} (denā) — to give}

\label{lesson:HI-C42-give}



\begin{quote}
Before the new one: what did \emph{khaṛā honā} mean?
\end{quote}

\subsection*{You'll want to know: \hi{देना}}
\textbf{\hi{देना}} (\emph{denā}) — \textquotedblleft{}to give\textquotedblright{}.

The plain 'give', and one of the light verbs Hindi uses to build compounds of its own.

\begin{cousinweb}
\textbf{\hi{देना}} (\emph{denā}) is a tadbhava, a word inherited from Sanskrit and reshaped by sound change. It comes from \emph{dadāti}, \textquotedblleft{}gives\textquotedblright{} (CDIAL 6141).
\end{cousinweb}

\begin{grammarlens}[title={What you've built}]
One more everyday action. That closes the run: four everyday actions, each met on its own.
\end{grammarlens}

\subsection*{Your turn}
Say these aloud:

\begin{itemize}
\raggedright
  \item \emph{denā}
  \item it again, to someone standing in a doorway
  \item it after \emph{khaṛā honā}, so the two sit together
\end{itemize}

\begin{tcolorbox}[breakable,colback=teal!4,colframe=teal!35!black,title={Before you move on}]
What does \hi{देना} mean? (\textquotedblleft{}to give\textquotedblright{}.) Say it once more.
\end{tcolorbox}
